\section{Related Work}
\label{sec:related_work}

\subsection{Deep Learning for CT Image Analysis}

Convolutional neural networks have rapidly become the dominant approach
for classification tasks on medical images.
Early work adapting ImageNet-pretrained architectures to CT slice
classification showed that transfer learning is effective despite the
domain gap between natural images and greyscale CT
slices~\cite{shin2016deeplearning,tajbakhsh2016convolutional}.
The key insight enabling effective transfer is the conversion of CT
Hounsfield Units to a pseudo-RGB representation via multiple windowing
presets---a technique that maps different tissue contrast ranges into
separate colour channels and has since become standard practice for CT
classification~\cite{chilamkurthy2018cq500,wang2021deeplearning}.

\subsection{Intracranial Hemorrhage Detection}

Chilamkurthy et al.~\cite{chilamkurthy2018cq500} were among the first to
demonstrate that deep learning at the study level can match or exceed
radiologist performance on several types of critical brain findings on
NCCT, releasing the CQ500 dataset used in this work.
Subsequent methods have explored diverse architectural choices, including
attention gates~\cite{schlemper2019attention}, squeeze-and-excitation
blocks~\cite{hu2018squeeze}, and multi-scale feature fusion
strategies~\cite{ye2019hemorrhage}.

Sage and Badura~\cite{sage2020doublebranch} proposed a double-branch
ResNet-50 framework that combines two complementary 3-channel RGB
representations of each CT slice: Branch~\#1 stacks the brain, subdural,
and bone windowed images of the same slice as three colour channels,
while Branch~\#2 stacks three consecutive slices $(s_{t-1}, s_t, s_{t+1})$
all in the subdural window to capture local spatial context.
Features from both branches are concatenated and passed to an external
classifier (SVM or Random Forest), with the random forest variant
achieving F1 scores of 89.1\%, 72.7\%, 93.1\%, 96.6\%, and 89.9\%
for SDH, EDH, IPH, IVH, and SAH respectively.
Their work highlights that multi-window intensity representations and
adjacent-slice spatial context are complementary rather than
mutually exclusive, motivating the use of the three-window approach
in our Stage-1 CNN preprocessing.

The 2019 RSNA Intracranial Hemorrhage Detection
challenge~\cite{rsna2019} marked a turning point.
Wang et al.~\cite{wang2021deeplearning} describe the winning two-stage
pipeline: an ensemble of three independently trained slice-level CNNs
(DenseNet-121, DenseNet-169, SE-ResNeXt-101) followed by a bidirectional
GRU that processes the per-slice logit sequences within each CT study.
This architecture directly inspired the present work.
Importantly, the original paper reports that study-level AUC improves by
1--3 percentage points when the sequence model is added, and by a further
margin when slice thickness is included as a contextual cue.

A subsequent systematic review by Cho et al.~\cite{cho2021systematic}
examined over 60 deep learning methods for ICH detection and classification,
concluding that (i) ensemble models consistently outperform single networks,
(ii) sequence modelling of scan-level context is beneficial but
underexplored in the literature, and (iii) few works report confidence
intervals, making cross-study comparisons unreliable.
Our work directly addresses points (ii) and (iii).

\subsection{Sequence Models for Medical Imaging}

Recurrent neural networks and their gated variants were first applied to
medical imaging for video-based echocardiography
analysis~\cite{ouyang2020video} and for predicting prognosis from
longitudinal imaging series.
For CT-based pathology, the key challenge is that a single CT study
comprises a variable-length ordered sequence of axial slices
(typically 20--60 for brain protocols), and the number of positive slices
is often small and concentrated around the lesion.
The LSTM-based sequence model of~\cite{ye2019hemorrhage} demonstrated
that aggregating slice-level features improves sensitivity for small
hemorrhages such as intraventricular bleeds.

Bidirectional GRUs have become a preferred recurrent layer because they
propagate context in both the superior-to-inferior and inferior-to-superior
directions simultaneously, which is important when the hemorrhage boundary
is ambiguous at slice boundaries~\cite{wang2021deeplearning}.
Our Stage-2 model follows this design and additionally conditions on slice
thickness to account for the fact that thicker slices cover more anatomy
per frame and therefore have different inter-slice correlation
characteristics.

\subsection{Ensemble Methods and Uncertainty Quantification}

Multi-model ensembling by probability averaging is a well-established
approach for reducing prediction variance in deep learning
classifiers~\cite{lakshminarayanan2017simple}.
In the medical imaging context, ensemble averaging of backbone CNNs trained
with different architectures provides complementary feature representations
and has been the dominant strategy in top-performing challenge
submissions~\cite{wang2021deeplearning,rsna2019}.

Uncertainty quantification through bootstrap confidence intervals is
recommended by clinical AI evaluation
guidelines~\cite{topol2019highperformance} and by reporting standards such
as TRIPOD~\cite{moons2015tripod}.
Despite this, the majority of published ICH detection papers report only
point-estimate AUC values~\cite{cho2021systematic}.
Our implementation computes 95\,\% bootstrap CIs for all reported AUC,
sensitivity, and specificity metrics.
